\section{Introduction}\label{sec:intro}

Federated analytics computes statistics across institutions without moving raw data to a central place~\cite{wang2022federated,elkordy2023federated}. A recent line of systems puts a large language model (LLM) agent in charge of planning such analyses: LAFA, which its authors describe as ``the first system that integrates LLM-agent-based data analytics with FA'', turns natural-language queries into optimised federated-analytics workflows~\cite{ji2025lafa}, agent frameworks such as LangGraph orchestrate stateful agents~\cite{langgraph2026docs}, and tool protocols such as the Model Context Protocol carry each tool message as an HTTP request~\cite{mcp2026spec}. In this paper we consider the tool-calling form of this design: an orchestrating agent turns an analyst's question into tool calls to data sites and aggregates their replies.

Encryption protects the content of each call and reply but does not hide its length~\cite{rfc9846}, and its timing is visible to any on-path observer. Traffic analysis has repeatedly shown that such metadata is enough to infer what a user is doing: which web page is loaded~\cite{panchenko2016website,sirinam2018deep}, what an AI assistant is answering~\cite{weiss2024what,mcdonald2025whisper}, which LLM agent a user is running and for which task~\cite{zhang2025agentprint,pouryousef2026benchmark}, and which model architecture a federated-learning client trains~\cite{shuvo2025flare,shuvo2026flint}. In agent-orchestrated federated analytics the same observer, placed on the links between the orchestrator and the sites, sees which sites are called and the size and timing of every exchange, which may reveal how much data each site processes and what kind of query runs. In cross-institution health or finance analytics, participation and data volume can be sensitive on their own.

Padding, batching and differentially private shaping can hide this metadata~\cite{cai2014systematic,wang2017walkie,sabzi2024netshaper}, but they are not free. Every padded byte is sent and received, every cover message wakes a node, and a shaper that releases data at fixed ticks keeps nodes active. On edge devices and in large server deployments the energy adds up, and the cost of anonymity traffic has already been shown to be material on smartphones~\cite{hugenroth2023powering}. Sustainable data infrastructure needs to know what metadata protection costs in joules, and operators need a way to spend a fixed energy budget where it reduces leakage most. Energy-constrained cover traffic has a long history in sensor-network location privacy~\cite{kamat2005enhancing,shao2008towards}, but our search found no work that studies orchestrator-to-site metadata in agentic federated analytics or evaluates its defences under an energy budget (\cref{sec:related}).

We present \emph{Green Shaping}. It shapes traffic at the tool-call layer, where whole calls and replies are visible before encryption, with three primitives: bucketed padding, synchronised rounds with cover sites, and a differentially private (DP) release of queued bytes at fixed ticks in the style of NetShaper~\cite{sabzi2024netshaper}. We formulate a controller that would choose the primitives' parameters from a calibrated per-node energy model and an attacker retrained on shaped traffic, subject to a joule budget per node. The completed implementation study evaluates correctness and observable resource costs; it does not validate a joule-constrained privacy optimum.

The contributions combine a scoped design and analysis with a reproducible engineering evaluation:
\begin{enumerate}[label=(\arabic*),leftmargin=*,noitemsep,topsep=2pt]
\item A system and threat model for passive inference of site participation, contribution size and query family from orchestrator-to-site metadata (\cref{sec:model}).
\item The Green Shaping design: three tool-call-layer primitives, the deployment conditions they need, a per-node energy cost model and a budget-constrained selection algorithm, with a reference implementation (\cref{sec:design}).
\item An analysis that scopes what each primitive can and cannot guarantee (\cref{sec:analysis}): a bound on the information a padded message size can carry; an energy floor at every site for hiding participation, so that the network-wide cost grows with the number of sites; a DP guarantee for the release mechanism under feedback between links, with a bound on any attacker's accuracy and an exact expression for its padding cost, which for the queue-length release grows with the square of the number of ticks; and the order in which primitives must be combined.
\item A bounded implementation study on public HIGGS data: TRAIN/SELECT calibration over 20 logical partitions, typed-plan and numerical-summary validation with two local LLMs, and 168 successful sessions through a four-instance EC2 deployment with captures at both sides of a common observer link (\cref{sec:eval,sec:results}). The study measures correctness, bytes, latency and recovery integrity, with retrospective question-disjoint metadata classification and independently validated fixed-arrival stress tests. It does not measure physical joules or edge-device performance, or establish a final-test privacy ranking.
\item A public-contract completion bound and prototype admission layer with per-node modelled budgets. Synthetic deployment and workflow-envelope studies compare conservative tick accounting, occupied-tick accounting, public clipping and constant cover. Paired loopback TLS acquisition additionally tests the envelope against two periodic-cover frequencies (\cref{sec:admission,sec:stress,sec:localtls}). Multi-host validation of these new candidates remains open.
\end{enumerate}

\Cref{sec:related} reviews related work. \Cref{sec:model,sec:design,sec:analysis} present the model, design and analysis. \Cref{sec:eval,sec:results} give the methodology and measured results, \cref{sec:discussion} discusses implications and limitations, and \cref{sec:conclusion} concludes.
